\section{Learning Signals by Credited Object}
\label{sec:learning}

\begin{figure}[t]
\centering
\resizebox{\columnwidth}{!}{%
\begin{tikzpicture}[font=\scriptsize,
  tok/.style={draw, minimum width=9mm, minimum height=6mm, align=center, font=\tiny},
  lvl/.style={font=\tiny, anchor=east}]
% one span of a trajectory: b_t (written belief) then a_t (action), then o_{t+1}
\node[tok, fill=yellow!25] (b1) at (0,0) {belief\\$\belief_t$};
\node[tok, fill=yellow!25, right=0mm of b1] (b2) {\op{Predict}};
\node[tok, fill=gray!15, right=0mm of b2] (r1) {step};
\node[tok, fill=gray!15, right=0mm of r1] (r2) {step};
\node[tok, fill=blue!12, right=0mm of r2] (a1) {action\\$a_t$};
\node[tok, dashed, right=1.5mm of a1] (o1) {$o_{t+1}$};
\node[tok, fill=gray!15, right=1.5mm of o1] (nx) {$\cdots$};
% credit brackets
\draw[decorate, decoration={brace, amplitude=2pt}, thick] ([yshift=4mm]b1.north west) -- ([yshift=4mm]a1.north east) node[midway, above=2pt, font=\tiny] {\textbf{output / trajectory}: reward the response or every step};
\draw[decorate, decoration={brace, amplitude=2pt, mirror}, thick] ([yshift=-4mm]a1.south west) -- ([yshift=-4mm]a1.south east) node[midway, below=2pt, font=\tiny] {\textbf{interaction}: reward the action};
\draw[decorate, decoration={brace, amplitude=2pt, mirror}, thick, red!60!black] ([yshift=-11mm]b1.south west) -- ([yshift=-11mm]b2.south east) node[midway, below=2pt, font=\tiny, text=red!60!black] {\textbf{state}: reward the belief, signed by $o_{t+1}$};
\draw[-{Latex}, red!60!black, thick] (o1.south) to[out=-90, in=-90, looseness=0.9] node[pos=0.5, below, font=\tiny, text=red!60!black] {surprise} ([xshift=2mm,yshift=-1mm]b2.south);
\end{tikzpicture}}
\caption{The credited object on one span of a trajectory. Output- and trajectory-level credit reaches text generated from the context; interaction-level credit reaches the action; state-level credit reaches the written belief, and its sign comes from the observation returned after the action.}
\label{fig:credit}
\end{figure}

\begin{table}[t]
\centering\scriptsize
\setlength{\tabcolsep}{3pt}
\renewcommand{\arraystretch}{1.08}
\resizebox{\columnwidth}{!}{%
\begin{tabular}{@{}llll@{}}
\toprule
Level & Credited object & Signed by $o_{t+1}$? & Representative \\
\midrule
Output & finished response & no & RLHF, DPO \\
Trajectory & sampled steps & no & PRMs, GRPO, LATS \\
Interaction & actions, tool calls & no (task reward) & ArCHer, RAGEN, HiPER \\
Interaction & memory operations & no (operation reward) & MEM1, Memento, Memory-R2 \\
State & written belief & closed-world oracle & CBM \\
State & claim + certainty & calibration target & Agent-BRACE \\
State & posterior tag & Bayesian teacher & BOND \\
State & written belief & \emph{own prediction} & (none published) \\
\bottomrule
\end{tabular}}
\caption{Learning signals by credited object and where the sign of the signal comes from (Section~\ref{sec:learning}); the full version with feedback source, granularity, and evidence level is Table~\ref{tab:learning} in the appendix.}
\label{tab:learningmain}
\end{table}


Existing surveys organise post-training by where the reward comes from, where the trajectories come from, and how dense the signal is \citep{wu2025rewards,zhang2026reasoning,zheng2025prm}. We hold those fixed and vary one thing, the object that receives the gradient (Figure~\ref{fig:credit}, Table~\ref{tab:learningmain}). The criterion is operational: a method is at the \emph{state level} if the quantity that receives a gradient is a claim about the environment and its sign depends on an observation returned after the action. A method that rewards a memory write regardless of whether the written claim was true is at the interaction level, however dense its credit.

\subsection{Output and Trajectory}
\label{sec:learn:traj}

Supervised fine-tuning and preference optimisation credit the completed response \citep{ouyang2022training,wei2022finetuned,rafailov2023direct}, and self-consistency and self-improvement select the responses to train on \citep{wang2022self,huang2022large}. Trajectory-level credit began by imitating and searching over agent trajectories \citep{chen2023fireact,zeng2023agenttuning,xi2024agentgym,zhou2023language,putta2024agent,fu2025agentrefine}, then credited sampled steps directly through PPO, GRPO, and process reward models \citep{schulman2017proximal,shao2024deepseekmath,lightman2023lets,zheng2025prm,zhang2026selfevolving,tan2026hindsight}.

\subsection{Interaction}
\label{sec:learn:inter}

Agentic reinforcement learning credits actions over several turns \citep{zhou2024archer,qi2024webrl,wang2025ragen,chen2025reinforcement,zhou2025sweetrl,feng2025group,dong2025agentic}, and the open question is how to make that credit denser without making it wrong, by hierarchical advantages, counterfactual estimators over shared states, potential-based shaping, or verifier-structured credit \citep{peng2026hiper,wang2026capo,diao2026hipif,wang2026groupgraph,ren2026iapo,liao2026tcpo,he2026turns,liu2026pica,xie2026tips,zheng2026pgpo,zhang2026mica,li2026vict,wang2026teach,guo2026hidifftir,qi2026stapo,zhao2026aem,lin2026skillc}; denser is not free \citep{modecrua2026multiturn,pu2026explore}. The nearest this level comes to the state is a group that credits the write to a memory, from a rewritten case bank to an RL-trained overwrite of a fixed written state \citep{zhou2025memento,zhou2025mem1,yu2025memagent,yan2026memoryr2,wang2026tamtrl,ning2026evoharness,wang2026slearl,hu2026collaborative}. By the criterion above these remain interaction-level: they credit the \emph{operation}, so a well-executed update to a wrong claim is still rewarded.

\subsection{State}
\label{sec:learn:state}

Three published methods meet the criterion. CBM trains with belief-state rewards from symbolic verification and cuts belief-management failures by about 71\% \citep{xu2026should}. Agent-BRACE trains belief model and policy jointly with verbalised certainty as the belief's output \citep{singh2026agent}. BOND distils a Bayesian teacher's posterior into an 8B student whose Brier score beats a 70B prompted baseline \citep{cui2026distilling}.

\insight{3}{\textbf{Credit has been getting denser, not deeper.} Agentic reinforcement learning has refined which action to prefer; no method changed what the action is conditioned on. The two axes are orthogonal, and the first system to combine them will need a sign for its belief that the environment, not an oracle, supplies.}
